\documentclass[11pt]{article}
\usepackage[a4paper,margin=1in]{geometry}
\usepackage{amsmath,amssymb,amsthm,mathtools}
\usepackage{booktabs}
\usepackage{enumitem}
\usepackage[hidelinks]{hyperref}
\usepackage{microtype}

\newtheorem{theorem}{Theorem}[section]
\newtheorem{proposition}[theorem]{Proposition}
\newtheorem{lemma}[theorem]{Lemma}
\newtheorem{corollary}[theorem]{Corollary}
\newtheorem{remark}[theorem]{Remark}
\newtheorem{definition}[theorem]{Definition}

\DeclareMathOperator{rad}{rad}
\DeclareMathOperator{Card}{Card}
\DeclareMathOperator{ord}{ord}

\title{Closure Gates and Countermodels in Multi-Route Research\\
on the $abc$ Conjecture}
\author{ChatGPT}
\date{28 August 2026}

\begin{document}
\maketitle

\begin{abstract}
We report a proof-audit program for the $abc$ conjecture in which every
proposed route is reduced to a small, explicitly quantified closure gate and
then separated from its already verified downstream consequences.  The
program covers a smooth-number counterexample construction, exceptional-set
amplification, inter-universal Teichm\"uller theory, Frey--Szpiro height
methods, and a Pell--Chebyshev reduction.  We prove five unconditional
results.  First, negligibility relative to an ambient interval does not imply
negligibility relative to a vanishing smooth-number main term.  Second, the
correct exponent condition in a near-prime-power construction is a linear
logarithmic-slope inequality, not a power inequality.  Third, a proposed
ambient high-$\omega$ tail estimate at the $x^{3/5}$ scale is contradicted by
multiples of primorial-type moduli, which already supply an $x^{4/5}$ lower
bound along an explicit subsequence.  Fourth, a power-saving exceptional-set
bound can yield finiteness only if a descendant construction crosses a precise
amplification threshold.  Fifth, neither an existentially good output of a
multivalued construction nor a nonnegative average of the existing Frey
models supplies the pointwise or coefficient-one conclusion required by
$abc$.  Lean 4 formalizations are provided for the scalar, finite-set, and
quantifier-theoretic cores.  No unconditional proof or disproof of the $abc$
conjecture is claimed.  The main contribution is an auditable map of the
remaining global obligations.
\end{abstract}

\section{The target and the audit discipline}

For coprime positive integers $a,b,c$ satisfying $a+b=c$, the classical
$abc$ conjecture asserts that for every $\varepsilon>0$ there is a constant
$C_\varepsilon>0$ such that
\begin{equation}\label{eq:abc}
 c\le C_\varepsilon\,\rad(abc)^{1+\varepsilon}.
\end{equation}
Equivalently,
\begin{equation}\label{eq:abclog}
 \log c\le (1+\varepsilon)\log\rad(abc)+O_\varepsilon(1).
\end{equation}
The repository accompanying this paper encodes the logarithmic form with a
uniform additive constant.  The central methodological rule is that a Lean
proof of a conditional implication is not a Lean proof of $abc$ unless every
premise is either constructed or imported from an independently accepted
unconditional theorem with matching quantifiers.

We use the following three-way classification.
\begin{enumerate}[label=(\roman*)]
\item A \emph{kernel theorem} is fully checked in Lean from the stated
      hypotheses.
\item A \emph{closure gate} is an explicit premise whose proof would complete
      a route.
\item A \emph{countermodel} disproves a proposed inference or intermediate
      assertion without deciding $abc$ itself.
\end{enumerate}
This distinction prevents an inhabited interface, an existential witness, or
an asymptotic with the wrong reference scale from being counted as the missing
arithmetic theorem.

\section{Smooth numbers: the relative extraction gate}

Let $h=h(x)$ be a short-interval length, let
\[
 S_x=\{n\in[x,x+h]:P^+(n)\le y\},
\]
and let $B_x\subseteq S_x$ be the subset whose number of distinct prime
factors is too large for radical control.  Suppose that a smooth-number theorem
gives
\begin{equation}\label{eq:smoothmain}
 |S_x|=h\rho(u)(1+o(1)),\qquad u=\frac{\log x}{\log y},
\end{equation}
where $\rho(u)\to0$.

The extraction problem is finite and exact: one needs $|B_x|<|S_x|$.  An
estimate $|B_x|=o(h)$ is not sufficient because the main set itself has
vanishing density in the interval.

\begin{proposition}[Ambient-scale countermodel]\label{prop:ambient}
There exist sequences of finite sets $B_n\subseteq S_n$ and ambient scales
$h_n\to\infty$ such that
\[
 |B_n|=o(h_n),\qquad B_n=S_n\quad\text{for every }n.
\]
Consequently ambient negligibility does not imply the existence of an element
of $S_n\setminus B_n$.
\end{proposition}

\begin{proof}
Take $h_n=n+1$ and $B_n=S_n=\{0\}$.  Then $|B_n|/h_n\to0$, but the complement
$S_n\setminus B_n$ is empty for every $n$.
\end{proof}

\begin{corollary}[Correct extraction scale]\label{cor:relative}
A sufficient asymptotic condition for extraction from \eqref{eq:smoothmain} is
\[
 |B_x|=o\!\left(h\rho(u)\right).
\]
More generally it suffices to prove the exact inequality $|B_x|<|S_x|$.
\end{corollary}

There is a second scale obstruction.  Dickman decay is faster than every
fixed negative power of $u$.  Thus an ambient error of size $h/u^A$ need not
be small relative to $h\rho(u)$; indeed
\[
 \frac{h/u^A}{h\rho(u)}=\frac{1}{u^A\rho(u)}\longrightarrow\infty.
\]
The Lean module \texttt{SmoothMainTermScaleAudit.lean} formalizes the finite
countermodel, the exact cardinality extraction lemma, and a
multiplication-only version of this super-polynomial obstruction.

\section{Near-prime powers: the correct logarithmic exponent}

Consider a prospective counterexample family
\begin{equation}\label{eq:nearpp}
 x=p^k,\qquad c=x+a,\qquad 0<a\le x^\theta,\qquad p\nmid c.
\end{equation}
Assume that
\begin{equation}\label{eq:c-rad-slope}
 \log\rad(c)\le (\eta+o(1))\log x.
\end{equation}
Then
\[
 \rad(ap^kc)\le a\,p\,\rad(c)
 \le x^{\theta+1/k+\eta+o(1)}.
\]
The exponent in \eqref{eq:abc} is therefore controlled by the linear
logarithmic condition
\begin{equation}\label{eq:correct-slope}
 (1+\varepsilon)\left(\theta+\frac1k+\eta\right)<1.
\end{equation}

\begin{proposition}[Power-condition counterexample]\label{prop:powerwrong}
The inequality
\[
 \left(\theta+\frac1k\right)^{1+\varepsilon}<1
\]
does not imply \eqref{eq:correct-slope}, even when $\eta=0$.
\end{proposition}

\begin{proof}
Take $\theta=3/5$, $k=5$, and $\varepsilon=1$.  Then
\[
 \left(\frac35+\frac15\right)^2=\frac{16}{25}<1,
 \qquad
 2\left(\frac35+\frac15\right)=\frac85>1.
\]
\end{proof}

For the interval exponent $\theta=3/5$, the exact corrected threshold is
\begin{equation}\label{eq:threefifths}
 (1+\varepsilon)\left(\frac35+\frac1k\right)<1
 \quad\Longleftrightarrow\quad
 5(1+\varepsilon)<k(2-3\varepsilon).
\end{equation}
Hence no positive $k$ works once $\varepsilon\ge2/3$, while for small fixed
$\varepsilon$ one may choose a sufficiently large $k$.  The radical loss
$\eta$ must fit the remaining budget
\[
 \eta<\frac{1}{1+\varepsilon}-\theta-\frac1k.
\]
These identities and the counterexample in Proposition~\ref{prop:powerwrong}
are formalized in \texttt{NearPrimePowerExponentAudit.lean}.

\begin{theorem}[Subcritical radical-slope disproof gate]\label{thm:slopegate}
Let $(a_n,b_n,c_n)$ be an unbounded family of positive pairwise coprime
solutions of $a_n+b_n=c_n$.  Suppose that for some $q,K\in\mathbb R$,
\[
 \log\rad(a_nb_nc_n)\le q\log\max\{a_n,b_n,c_n\}+K
\]
for every $n$, and that $(1+\varepsilon)q<1$ for some $\varepsilon>0$.
Then the $abc$ conjecture is false.
\end{theorem}

\begin{proof}
If \eqref{eq:abclog} held with additive constant $C_\varepsilon$, then
\[
 \left(1-(1+\varepsilon)q\right)
 \log\max\{a_n,b_n,c_n\}
 \le (1+\varepsilon)K+C_\varepsilon.
\]
The coefficient on the left is positive, contradicting unbounded height.
\end{proof}

Theorem~\ref{thm:slopegate} is formalized in
\texttt{SubcriticalRadicalSlopeDisproofGate.lean}.  It is a complete final
implication, not a construction of the required family.

\section{An elementary refutation of an ambient high-$\omega$ hypothesis}

A proposed smooth-number counterexample argument additionally uses an ambient
hypothesis of the form
\begin{equation}\label{eq:ambientomega}
 \#\{n\le x:\omega(n)>w(x)\}=o(x^{3/5}),
\end{equation}
with $w(x)=O(\log\log x)$.  This is false by a direct divisibility argument.

\begin{theorem}[Primorial-multiple lower bound]\label{thm:omega}
Let $S$ be a finite set of $r$ distinct primes and let
$Q=\prod_{p\in S}p$.  Then
\[
 \#\{n\le Q^5:\omega(n)\ge r\}\ge Q^4.
\]
Consequently, for any unbounded sequence of such products satisfying
$w(Q^5)<r$, the left side of \eqref{eq:ambientomega} is at least
$x^{4/5}$ along $x=Q^5$.
\end{theorem}

\begin{proof}
The numbers $Q,2Q,\ldots,Q^4Q$ are pairwise distinct and at most $Q^5$.
Every one is divisible by all primes in $S$, hence has at least $r$ distinct
prime factors.
\end{proof}

For the product $Q_r$ of the first $r$ primes, a crude iterative Bertrand
bound gives $\log\log(Q_r^5)=O(\log r)$, so every threshold
$w(x)=O(\log\log x)$ satisfies $w(Q_r^5)<r$ eventually.  Thus
\[
 \#\{n\le Q_r^5:\omega(n)>w(Q_r^5)\}
 \ge Q_r^4=(Q_r^5)^{4/5},
\]
and the ratio to $(Q_r^5)^{3/5}$ tends to infinity.  The finite counting core
is formalized in \texttt{HighOmegaAmbientTailRefutation.lean}.  This refutes
an intermediate hypothesis, not the possibility of a different conditional
distribution theorem inside the smooth set.

\section{Exceptional sets and the amplification threshold}

Suppose an exceptional-set theorem gives
\begin{equation}\label{eq:exceptional}
 E(X)\ll X^\beta.
\end{equation}
Such a power saving is compatible with an infinite lacunary set.  A proof of
finiteness needs a propagation theorem that turns each bad source object into
many distinct bad descendants.

\begin{theorem}[Amplification threshold]\label{thm:amplification}
Assume that a bad object of source height $H$ produces at least
$H^{\gamma-o(1)}$ pairwise distinct bad descendants, all of height at most
$H^{d+o(1)}$.  If \eqref{eq:exceptional} holds uniformly, then necessarily
\[
 \gamma\le d\beta.
\]
In particular, a construction with $\gamma>d\beta$ contradicts the
exceptional-set envelope along an unbounded sequence.
\end{theorem}

\begin{proof}
The descendants give
$H^{\gamma-o(1)}\le E(H^{d+o(1)})\ll H^{d\beta+o(1)}$.
Comparison of logarithmic exponents yields the claim.
\end{proof}

For the exponent $\beta=56/85$, quadratic height dilation demands
$\gamma>112/85$, and quartic dilation demands $\gamma>224/85$.  The exact
scalar contradiction is formalized in
\texttt{ExceptionalSetAmplificationThreshold.lean}.  The remaining arithmetic
problem is to construct descendants, prove bounded fibres after removing
symmetries and common factors, and retain one fixed quality exponent.

\section{IUT: existential output versus input-determined output}

Let $\mathcal P(x)$ be a set of possible outputs of a multivalued construction
and let $G(x,y)$ be the estimate required of an output.  The statement
\begin{equation}\label{eq:exists-output}
 \forall x\ \exists y\in\mathcal P(x),\quad G(x,y)
\end{equation}
does not imply that a distinguished output $\operatorname{sel}(x)$ satisfies
$G$.

\begin{proposition}[Two-output countermodel]\label{prop:output}
There is a one-input, two-output model in which \eqref{eq:exists-output} holds,
but the selected possible output is not good.
\end{proposition}

\begin{proof}
Declare both Boolean values possible, declare only \texttt{true} good, and
select \texttt{false}.
\end{proof}

Two sufficient repairs are immediate.
\begin{enumerate}[label=(\alph*)]
\item \emph{Uniqueness:} all possible outputs for one input are equal.
\item \emph{Fibre invariance:} if $y,z\in\mathcal P(x)$ and $G(x,y)$, then
      $G(x,z)$.
\end{enumerate}
The countermodel and both repair theorems are formalized in
\texttt{IUTOutputIdentificationAudit.lean}.  For the actual IUT route, the
closure gate is source-derived uniqueness, invariance under all relevant
indeterminacies, or a canonical quotient on which the normalized volume
estimate is well defined.  Merely including a theta-pilot region and its shell
containment as fields of an input structure does not prove that the actual
algorithm produces the required pointwise output.

\section{Frey--Szpiro: a convex-slope obstruction}

The existing explicit Frey models control the abc height through
discriminant-growth exponents $4$ and $5$, while a critical Szpiro inequality
has conductor slope $6$.  Suppose nonnegative weights $w_4,w_5$ are used to
combine the two estimates.  The resulting scalar ratio is
\begin{equation}\label{eq:frey-ratio}
 R(w_4,w_5)=\frac{6(w_4+w_5)}{4w_4+5w_5}.
\end{equation}

\begin{theorem}[Convex-slope barrier]\label{thm:frey}
For $w_4,w_5\ge0$ and $w_4+w_5>0$,
\[
 R(w_4,w_5)\ge\frac65>1.
\]
Therefore nonnegative averaging of the two existing discriminant estimates
cannot reach the coefficient one required by $abc$.
\end{theorem}

\begin{proof}
The inequality $R\ge6/5$ is equivalent, after multiplication by the positive
denominator, to
$30(w_4+w_5)\ge24w_4+30w_5$, i.e. $6w_4\ge0$.
\end{proof}

More generally, a nonnegative combination cannot improve the best ratio
$s_i/r_i$ among its components.  The modified integral height
$\log\max\{|c_4|^3,|\Delta|\}$ is different because it grows as
$6\log c+O(1)$, yielding the critical scalar ratio $6/6=1$.  The unresolved
gate is a uniform, non-circular global estimate for that exponent-six height.
Theorem~\ref{thm:frey} and the general two-model lemma are formalized in
\texttt{FreyConvexSlopeBarrier.lean}.

\section{Pell--Chebyshev: what a uniform theorem must supply}

The fixed-prime computations through $p=31$ leave a residual prime-index
family beginning at $p\ge37$.  Its exact scalar reduction has the form
\begin{equation}\label{eq:pell}
 (3AB+1)^{37}\le Z^2<(b+2)^4.
\end{equation}
Thus a pointwise lower bound just beyond the scale
$AB\asymp b^{4/37}$ would close all remaining prime indices after a finite
calculation.  Existing recurrence-support estimates control support rather
than valuation parity and, after conversion along the Pell orbit, remain at a
$b^{o(1)}$ scale.  The decisive theorem must therefore provide a fixed
positive-power lower bound for the relevant parity kernel, or an alternative
uniform Selmer/local obstruction independent of the prime index.  Continuing
with one class-group and Chabauty certificate for each prime cannot by itself
prove the infinite-index statement.

\section{The remaining closure matrix}

The current proof state may be summarized as follows.

\begin{center}
\begin{tabular}{p{0.18\textwidth}p{0.32\textwidth}p{0.38\textwidth}}
\toprule
Route & Kernel-checked or rigorously proved part & Remaining global obligation \\
\midrule
Smooth counterexample & Relative extraction scale; corrected exponent;
subcritical slope implies $\neg abc$ & Construct low-radical smooth integers
in the required short intervals on an unbounded prime-power subsequence \\
Exceptional set & Exact threshold $\gamma>d\beta$ & Produce enough distinct
quality-preserving descendants with controlled height and bounded fibres \\
IUT & Existential/selected-output countermodel; uniqueness and invariance
repairs & Derive the actual output and its volume estimate from genuine
Hodge-theater/Kummer data with correct quantifiers \\
Frey--Szpiro & Explicit local arithmetic; convex averaging barrier; conditional
modified-height gate & Prove a uniform exponent-six global conductor/height
estimate without importing an equivalent of $abc$ \\
Pell--Chebyshev & Fixed-index certificates and the exact $4/37$ residual &
Uniform positive-power parity-kernel bound or prime-independent
Selmer/local obstruction \\
\bottomrule
\end{tabular}
\end{center}

Each remaining item is a global theorem rather than a missing line of scalar
algebra.  Conversely, proving any one of the appropriate proof gates, or
constructing the smooth-number family in the disproof gate, would activate an
already explicit downstream chain.

\section{Formalization and trust boundary}

The Lean development uses Lean~4 and Mathlib.  The new modules associated with
this paper are:
\begin{verbatim}
SmoothMainTermScaleAudit.lean
SubcriticalRadicalSlopeDisproofGate.lean
NearPrimePowerExponentAudit.lean
HighOmegaAmbientTailRefutation.lean
ExceptionalSetAmplificationThreshold.lean
IUTOutputIdentificationAudit.lean
FreyConvexSlopeBarrier.lean
\end{verbatim}
The files contain no \texttt{sorry}, \texttt{admit}, or newly postulated
arithmetic axiom.  Their scope is intentionally limited to finite-set
arguments, exact rational inequalities, generic logarithmic implications, and
quantifier logic.  Analytic smooth-number estimates, IUT comparison theorems,
modified Szpiro, external Chabauty certificates, and any uniform parity-kernel
bound remain outside the kernel unless separately formalized and audited.

Passing Lean elaboration establishes that a theorem follows from its encoded
premises; it does not establish that an abstract structure has the intended
mathematical semantics.  For this reason every bridge interface must be
accompanied by an inhabitation audit and a source-construction theorem.

\section{Conclusion}

The multi-route program has not produced an unconditional proof or disproof of
the $abc$ conjecture.  It has, however, replaced several ambiguous gaps by
sharp mathematical obligations and eliminated three tempting but invalid
shortcuts: ambient-scale extraction for a vanishing main term, the wrong
near-prime-power exponent condition, and convex averaging of subcritical Frey
height models.  It also supplies an elementary contradiction to a proposed
ambient high-$\omega$ tail estimate and isolates the exact amplification and
output-identification thresholds required by two other routes.

The highest-priority positive targets are now: a source-derived IUT output
invariance theorem, a non-circular exponent-six modified-height inequality,
and a fixed positive-power parity-kernel estimate in the Pell--Chebyshev
family.  The highest-priority disproof target is a direct construction of
short-interval smooth integers whose radical slope fits
\eqref{eq:correct-slope}.  Only after one of these global gates is closed and
the complete dependency graph passes kernel and circularity audits should a
theorem named \texttt{abc\_conjecture} or its negation be merged into the
\texttt{main} branch.

\begin{thebibliography}{99}

\bibitem{Baker}
A.~Baker,
\emph{Logarithmic forms and the $abc$ conjecture},
various foundational works on effective Diophantine approximation.

\bibitem{BLT}
T.~D. Browning, J.~D. Lichtman, and J.~Ter\"av\"ainen,
work on power-saving bounds for the exceptional set in the $abc$ conjecture.

\bibitem{Carella2026}
N.~A. Carella,
\emph{Note on the Exceptional Set in the ABC Conjecture},
arXiv:2608.16764v2, 2026.

\bibitem{Dickman}
K.~Dickman,
\emph{On the frequency of numbers containing prime factors of a certain
relative magnitude}, Arkiv f\"or Matematik, Astronomi och Fysik, 1930.

\bibitem{Li2025}
R.~Li,
\emph{On the exceptional set in the $abc$ conjecture},
arXiv:2507.02885, 2025.

\bibitem{Lichtman2025}
J.~D. Lichtman,
\emph{The $abc$ conjecture is true almost always},
arXiv:2505.13991, 2025.

\bibitem{MochizukiIUT}
S.~Mochizuki,
\emph{Inter-universal Teichm\"uller theory I--IV},
Publ. Res. Inst. Math. Sci. 57 (2021).

\bibitem{ProjectLANA2026}
Project LANA,
\emph{Formal verification of inter-universal Teichm\"uller theory: interim
report}, 2026.

\end{thebibliography}

\end{document}
